\maketitle

\begin{abstract}
사진 몇 장으로 3D 장면을 재구성하는 방법은 크게 두 갈래로 나뉜다. 하나는 미리 학습한 모델에 사진을 넣으면 즉시 3D Gaussian을 출력하는 feed-forward 방식이고, 다른 하나는 새 장면마다 그 장면 전용으로 반복 최적화하는 per-scene optimization 방식이다.

본 연구는 새로운 알고리즘을 제안하는 대신, \textbf{입력 view 수, view 간 겹침 정도(co-visibility overlap), 사용 가능한 계산 시간 예산}을 통제 축으로 삼아 대표 시스템들(feed-forward: MVSplat~\cite{chen2024mvsplat}, DepthSplat~\cite{xu2025depthsplat}; per-scene optimization: 3D Gaussian Splatting~\cite{kerbl20233d}, FSGS~\cite{zhu2024fsgs})의 실용적 우위 지도(regime map)를 그린다.

동일한 feed-forward 초기값에서 표준 3DGS refinement를 켜고 끄는 통제 실험으로 최적화의 순효과를 분리하고, 카메라 기하학 기반 불확실도 분석과 통제된 depth noise 개입으로 우위가 갈리는 원인을 설명하며, 마지막으로 조건별 선택 가이드라인을 제시한다. \todo{실험 완료 후 핵심 수치 1--2개를 요약에 추가}
\end{abstract}

\noindent\todo{본 문서는 실험이 진행 중인 상태의 초안이다. 표/수치가 비어 있는 절은 명시적으로 표시했다. 산문(서론, 관련 연구, 프로토콜, 일부 분석)은 지금까지 확정된 내용을 담았다.}
